%*************************************
% บทที่ 3
%*************************************
\chapterTitle{3}{วิธีการดำเนินงานและการออกแบบระบบ}

% Project Methodology
\section{ขั้นตอนการดำเนินงาน}

\hspace*{1.5em} %ย่อหน้า
การพัฒนาแพลตฟอร์มคลัสเตอร์เสมือนสำหรับการศึกษามีขั้นตอนการดำเนินงานที่เป็นระบบ โดยแบ่งออกเป็น 5 ขั้นตอนหลัก ดังนี้

\begin{mycustomenum2}
    \item \textbf{การศึกษาและวิเคราะห์ความต้องการ} การสำรวจความต้องการของผู้ใช้งาน การศึกษาระบบที่มีอยู่เดิม และการกำหนดข้อกำหนดของระบบ
    \item \textbf{การออกแบบสถาปัตยกรรมระบบ} การออกแบบโครงสร้างระบบทั้งหมด การเลือกเทคโนโลยีที่เหมาะสม และการออกแบบฐานข้อมูล
    \item \textbf{การพัฒนาระบบ} การพัฒนาส่วนต่าง ๆ ของระบบตามสถาปัตยกรรมที่ออกแบบไว้
    \item \textbf{การทดสอบระบบ} การทดสอบการทำงานของระบบ การทดสอบประสิทธิภาพ และการทดสอบความปลอดภัย
    \item \textbf{การประเมินผลและปรับปรุง} การประเมินผลการใช้งานจริง และการปรับปรุงระบบตามผลการประเมิน
\end{mycustomenum2}

% Research and Requirements Analysis
\section{การศึกษาและวิเคราะห์ความต้องการ}

% 3.2.1
\subsection{การสำรวจความต้องการผู้ใช้งาน}

\hspace*{1.5em} %ย่อหน้า
ผู้จัดทำได้ดำเนินการสำรวจความต้องการจากกลุ่มผู้ใช้งานในภาควิชาเทคโนโลยีสารสนเทศ ได้แก่

\begin{mycustomitem}
    \item \textbf{อาจารย์} ต้องการระบบจัดการเครื่องเสมือนที่ใช้งานง่าย สำหรับการสอนและการวิจัย
    \item \textbf{นักศึกษา} ต้องการเข้าถึงเครื่องเสมือนสำหรับการเรียนและทำโครงงาน
    \item \textbf{ผู้ดูแลระบบ} ต้องการเครื่องมือติดตามและจัดการระบบที่มีประสิทธิภาพ
\end{mycustomitem}

% 3.2.2
\subsection{การวิเคราะห์ระบบเดิม}

\hspace*{1.5em} %ย่อหน้า
จากการศึกษาระบบที่มีอยู่เดิม พบปัญหาและข้อจำกัดดังนี้

\begin{mycustomitem}
    \item \textbf{การใช้งานที่ซับซ้อน} ต้องใช้ความรู้ทางเทคนิคสูงในการจัดการเครื่องเสมือน
    \item \textbf{ขาดระบบจัดการผู้ใช้} ไม่มีการแยกสิทธิ์การใช้งานตามบทบาท
    \item \textbf{ประสิทธิภาพต่ำ} การสร้างและลบเครื่องเสมือนใช้เวลานาน
    \item \textbf{ขาดการติดตาม} ไม่มีระบบ Monitoring การใช้ทรัพยากร
\end{mycustomitem}

% 3.2.3
\subsection{ข้อกำหนดของระบบ}

\hspace*{1.5em} %ย่อหน้า
จากการวิเคราะห์ความต้องการ ได้กำหนดข้อกำหนดของระบบดังนี้

\hspace*{1.5em} %ย่อหน้า
\textbf{ข้อกำหนดเชิงหน้าที่ (Functional Requirements)}
\begin{mycustomitem}
    \item ระบบต้องสามารถสร้าง แก้ไข และลบเครื่องเสมือนได้
    \item ระบบต้องรองรับการจัดการผู้ใช้แบบ Role-Based Access Control
    \item ระบบต้องมีระบบยืนยันตัวตนผ่าน OAuth 2.0
    \item ระบบต้องสามารถติดตามการใช้ทรัพยากรแบบ Real-time
    \item ระบบต้องส่งการแจ้งเตือนเมื่อมีการเปลี่ยนแปลงสถานะ
\end{mycustomitem}

\hspace*{1.5em} %ย่อหน้า
\textbf{ข้อกำหนดที่ไม่เกี่ยวกับหน้าที่ (Non-Functional Requirements)}
\begin{mycustomitem}
    \item \textbf{ประสิทธิภาพ} การสร้างเครื่องเสมือนต้องใช้เวลาไม่เกิน 3 นาที
    \item \textbf{ความพร้อมใช้งาน} ระบบต้องพร้อมใช้งานอย่างน้อย 99.5% ของเวลา
    \item \textbf{ความปลอดภัย} ข้อมูลต้องได้รับการเข้ารหัสและจัดเก็บอย่างปลอดภัย
    \item \textbf{ความยืดหยุ่น} ระบบต้องรองรับการขยายทรัพยากรได้ในอนาคต
\end{mycustomitem}

% System Architecture Design
\section{การออกแบบสถาปัตยกรรมระบบ}

% 3.3.1
\subsection{สถาปัตยกรรมระบบโดยรวม}

\begin{figure}[htbp]
    \centering
    \includegraphics[width=1.0\textwidth]{Image/web/context_diagram.png}
    \caption{แผนภาพบริบทของระบบแพลตฟอร์มคลัสเตอร์เสมือน}
    \label{fig:context_diagram}
\end{figure}

\hspace*{1.5em} %ย่อหน้า
ระบบแพลตฟอร์มคลัสเตอร์เสมือนได้รับการออกแบบในรูปแบบ Microservices Architecture โดยแบ่งเป็นส่วนประกอบหลัก 3 ส่วน คือ

\begin{mycustomitem}
    \item \textbf{Frontend} เว็บแอปพลิเคชันสำหรับผู้ใช้งาน พัฒนาด้วย Next.js และ React
    \item \textbf{Backend} เซิร์ฟเวอร์ API สำหรับการจัดการข้อมูลและการยืนยันตัวตน พัฒนาด้วย Elysia.js
    \item \textbf{Instance Management} เซอร์วิสจัดการเครื่องเสมือนและการสื่อสารกับ Proxmox VE
\end{mycustomitem}

\hspace*{1.5em} %ย่อหน้า
จากภาพที่ \ref{fig:context_diagram} แผนภาพบริบทข้างต้นแสดงให้เห็นถึงการเชื่อมต่อและการสื่อสารระหว่างส่วนประกอบต่าง ๆ ของระบบ โดยผู้ใช้งานสามารถเข้าถึงระบบผ่านเว็บแอปพลิเคชัน ที่จะสื่อสารกับ Backend API และระบบจัดการเครื่องเสมือนเพื่อให้บริการที่สมบูรณ์

% 3.3.2
\subsection{การออกแบบฐานข้อมูล}

\hspace*{1.5em} %ย่อหน้า
ระบบใช้ PostgreSQL เป็นฐานข้อมูลหลัก โดยออกแบบตารางหลักดังนี้

\begin{mycustomitem}
    \item \textbf{User Table} เก็บข้อมูลผู้ใช้งาน บทบาท และสิทธิ์การเข้าถึง
    \item \textbf{Instance Table} เก็บข้อมูลเครื่องเสมือน การกำหนดค่า และสถานะ
    \item \textbf{Template Table} เก็บข้อมูล Template สำหรับการสร้างเครื่องเสมือน
    \item \textbf{Activity Log Table} เก็บประวัติการใช้งานและการเปลี่ยนแปลง
    \item \textbf{Resource Usage Table} เก็บข้อมูลการใช้ทรัพยากรแบบ Time-series
\end{mycustomitem}

% 3.3.3
\subsection{การออกแบบ API และการสื่อสาร}

\hspace*{1.5em} %ย่อหน้า
ระบบใช้ RESTful API เป็นหลักในการสื่อสารระหว่างส่วนประกอบต่าง ๆ โดยมีการออกแบบดังนี้

\begin{mycustomitem}
    \item \textbf{Authentication Endpoints} สำหรับการ Login, Logout และ Token Management
    \item \textbf{Instance Management Endpoints} สำหรับการจัดการเครื่องเสมือน (CRUD Operations)
    \item \textbf{Monitoring Endpoints} สำหรับการติดตามการใช้ทรัพยากร
\end{mycustomitem}

% 3.3.4
\subsection{การออกแบบระบบ Message Queue}

\hspace*{1.5em} %ย่อหน้า
เพื่อรองรับการประมวลผลแบบ Asynchronous ระบบใช้ RabbitMQ เป็น Message Broker โดยออกแบบ Queue หลักดังนี้

\begin{mycustomitem}
    \item \textbf{Instance Creation Queue} สำหรับงานการสร้างเครื่องเสมือน
    \item \textbf{Instance Deletion Queue} สำหรับงานการลบเครื่องเสมือน
\end{mycustomitem}

% Tools and Technologies Used
\section{เครื่องมือและเทคโนโลยีที่ใช้ในการพัฒนา}

% 3.4.1
\subsection{เทคโนโลยี Frontend}

\hspace*{1.5em} %ย่อหน้า
\textbf{Next.js 14} เป็น React Framework ที่เลือกใช้เนื่องจากคุณสมบัติดังนี้
\begin{mycustomitem}
    \item \textbf{App Router} สำหรับการจัดการ Routing แบบใหม่
    \item \textbf{Server Components} ลดปริมาณ JavaScript ที่ส่งไปยัง Client
    \item \textbf{Static Site Generation} เพิ่มประสิทธิภาพการโหลดหน้าเว็บ
    \item \textbf{Built-in TypeScript Support} รองรับ Type Safety
\end{mycustomitem}

\hspace*{1.5em} %ย่อหน้า
\textbf{Tailwind CSS} ใช้สำหรับการจัดการ Styling เนื่องจาก
\begin{mycustomitem}
    \item \textbf{Utility-first} เขียน CSS ได้เร็วและมีประสิทธิภาพ
    \item \textbf{Responsive Design} รองรับการแสดงผลบนอุปกรณ์ต่าง ๆ
\end{mycustomitem}

% 3.4.2
\subsection{เทคโนโลยี Backend}

\hspace*{1.5em} %ย่อหน้า
\textbf{Elysia.js} เป็น Web Framework ที่ใช้กับ Bun Runtime โดยมีเหตุผลในการเลือกใช้ดังนี้
\begin{mycustomitem}
    \item \textbf{High Performance} มีประสิทธิภาพสูงกว่า Node.js Framework ทั่วไป
    \item \textbf{Type Safety} รองรับ TypeScript แบบเต็มรูปแบบ
    \item \textbf{OpenAPI Integration} สร้าง API Documentation อัตโนมัติ
    \item \textbf{Plugin Ecosystem} มี Plugin หลากหลายสำหรับการพัฒนา
\end{mycustomitem}

\hspace*{1.5em} %ย่อหน้า
\textbf{Prisma ORM} ใช้สำหรับการจัดการฐานข้อมูล เนื่องจาก
\begin{mycustomitem}
    \item \textbf{Type-safe Database Client} ป้องกันข้อผิดพลาดในการเขียน Query
    \item \textbf{Database Migrations} จัดการการเปลี่ยนแปลง Schema อย่างเป็นระบบ
    \item \textbf{Database Introspection} สร้าง Schema จากฐานข้อมูลที่มีอยู่
\end{mycustomitem}

% 3.4.3
\subsection{Infrastructure และ DevOps}

\hspace*{1.5em} %ย่อหน้า
\textbf{Docker} ใช้สำหรับการ Containerization ของแอปพลิเคชัน เพื่อ
\begin{mycustomitem}
    \item \textbf{Environment Consistency} สภาพแวดล้อมการพัฒนาและ Production เหมือนกัน
    \item \textbf{Easy Deployment} ง่ายต่อการ Deploy และ Scale
    \item \textbf{Resource Isolation} แยกทรัพยากรระหว่างแอปพลิเคชัน
\end{mycustomitem}

\hspace*{1.5em} %ย่อหน้า
\textbf{Docker Compose} ใช้สำหรับการจัดการ Multi-container Application
\begin{mycustomitem}
    \item \textbf{Service Orchestration} จัดการการเริ่มต้นและหยุดทำงานของ Service
    \item \textbf{Network Management} จัดการเครือข่ายระหว่าง Container
    \item \textbf{Volume Management} จัดการการแชร์ข้อมูลระหว่าง Container
\end{mycustomitem}

% 3.4.4
\subsection{เครื่องมือการพัฒนา}

\hspace*{1.5em} %ย่อหน้า
\textbf{Version Control และ Collaboration}
\begin{mycustomitem}
    \item \textbf{Git} สำหรับการควบคุมเวอร์ชันของโค้ด
    \item \textbf{GitHub} สำหรับการเก็บโค้ดและการทำงานร่วมกัน
    \item \textbf{GitHub Actions} สำหรับ CI/CD Pipeline
\end{mycustomitem}

\hspace*{1.5em} %ย่อหน้า
\textbf{Code Quality และ Testing}
\begin{mycustomitem}
    \item \textbf{ESLint และ Prettier} สำหรับการตรวจสอบและจัดรูปแบบโค้ด
    \item \textbf{TypeScript} สำหรับ Static Type Checking
    \item \textbf{Jest} สำหรับ Unit Testing
\end{mycustomitem}

% Planning and Project Management
\section{การวางแผนการพัฒนาและการจัดการโครงการ}

% 3.5.1
\subsection{วิธีการพัฒนา (Development Methodology)}

\hspace*{1.5em} %ย่อหน้า
โครงงานนี้ใช้วิธีการพัฒนาแบบ Agile Development โดยแบ่งการพัฒนาเป็น Sprint ขนาด 2 สัปดาห์ แต่ละ Sprint จะมีขั้นตอนดังนี้

\begin{mycustomenum2}
    \item \textbf{Sprint Planning} วางแผนงานที่จะทำใน Sprint นั้น ๆ
    \item \textbf{Daily Standup} ติดตามความคืบหน้าและปัญหาที่เกิดขึ้น  
    \item \textbf{Development} การพัฒนาตาม User Stories ที่กำหนด
    \item \textbf{Testing} การทดสอบฟีเจอร์ที่พัฒนาเสร็จ
    \item \textbf{Sprint Review} ทบทวนผลงานและรับ Feedback
    \item \textbf{Sprint Retrospective} หาจุดปรับปรุงสำหรับ Sprint ถัดไป
\end{mycustomenum2}

% 3.5.2
\subsection{แผนการพัฒนาตามระยะเวลา}

\hspace*{1.5em} %ย่อหน้า
การพัฒนาแบ่งออกเป็น 4 Phase หลัก ตามลำดับดังนี้

\hspace*{1.5em} %ย่อหน้า
\textbf{Phase 1: Foundation (สัปดาห์ที่ 1-4)}
\begin{mycustomitem}
    \item ตั้งค่า Development Environment
    \item พัฒนาระบบ Authentication และ User Management
    \item สร้าง Database Schema และ API พื้นฐาน
    \item พัฒนา UI Components พื้นฐาน
\end{mycustomitem}

\hspace*{1.5em} %ย่อหน้า
\textbf{Phase 2: Core Features (สัปดาห์ที่ 5-8)}
\begin{mycustomitem}
    \item พัฒนาระบบจัดการเครื่องเสมือนพื้นฐาน (Create, Read, Update, Delete)
    \item เชื่อมต่อกับ Proxmox VE API
    \item พัฒนา Message Queue System
    \item สร้าง Dashboard หลักสำหรับผู้ใช้งาน
\end{mycustomitem}

\hspace*{1.5em} %ย่อหน้า
\textbf{Phase 3: Advanced Features (สัปดาห์ที่ 9-12)}
\begin{mycustomitem}
    \item พัฒนาระบบ Monitoring และ Resource Tracking
    \item เพิ่มระบบ Notification
    \item พัฒนา Advanced VM Management Features
    \item เพิ่มระบบ Template Management
\end{mycustomitem}

\hspace*{1.5em} %ย่อหน้า
\textbf{Phase 4: Testing และ Optimization (สัปดาห์ที่ 13-16)}
\begin{mycustomitem}
    \item Integration Testing และ Performance Testing
    \item Security Testing และ Load Testing  
    \item UI/UX Improvements
    \item Documentation และ Deployment Preparation
\end{mycustomitem}
